% II. АНАЛИЗ НА ВЪЗМОЖНИТЕ РЕШЕНИЯ И ОБОСНОВКА НА ИЗБОРА
% Съпоставка по четири оси: как се чете C++, в какво се извежда диаграмата, как се
% описва съответствието клас към таблица и с какви външни зависимости се строи.
\section{Анализ на възможните решения и обосновка на избора}
\label{sec:alternatives}

\subsection{Прочитане на изходния текст на C++}

Изборът на начин за прочитане на изходния текст предопределя горната граница на всичко
останало. Разгледани са четири възможности. Първата е анализ с регулярни изрази върху
текста на файловете, който се пише за часове, но се проваля при шаблони, пространства
от имена, макроси и всяко условно компилиране, тоест точно там, където реалният код е
интересен. Втората е използване на съществуващ инструмент за документация, който вече
разбира езика, и извличане на модела от неговия изход. Третата е свързване с
libclang~\cite{libclang_docs}, тоест със самия компилаторен фронтенд през програмен
интерфейс на C. Четвъртата е \term{изпълнение на компилатора като външен процес} с
ключа \code{-Xclang -ast-dump=json} и разбор на изведеното от него абстрактно
синтактично дърво в JSON~\cite{clang_ast_json,rfc8259}.

Първите две отпадат по същество. Текстовият анализ няма достъп до разрешените имена и
следователно не може да отговори на въпроса дали \code{Wheel} в един файл и
\code{Wheel} в друг са един и същ клас. Doxygen с Graphviz~\cite{doxygen_manual,%
ellson2001graphviz} е най-близкият утвърден инструмент и наистина чете C++ и рисува
диаграми на наследяване и на сътрудничество. Три обстоятелства обаче го правят негоден
за целта тук. Диаграмите му са изображения, предназначени за четене от човек, а не
модел, върху който може да се разсъждава машинно. Той не различава разновидностите на
асоциацията по семантика на собствеността, а само отбелязва наличието на връзка. И
най-важното, той няма права посока: от неговия изход не се поражда нито код, нито схема.
Тоест той решава едната половина от задачата и по устройство не може да реши другата.

Изборът между третата и четвъртата възможност е същинският. И двете използват един и
същ фронтенд и дават един и същ семантичен резултат: разрешени имена, типове след
заместване на шаблонните аргументи и точно местоположение на всяка декларация. Разликата
е изцяло в начина на свързване и е обобщена в Табл.~\ref{tab:frontend}.

\begin{table}[H]
\caption{Съпоставка на начините за достъп до компилаторния фронтенд}
\label{tab:frontend}
\centering
\fontsize{11}{13}\selectfont
\begin{tabular}{@{}L{3.6cm}L{4.6cm}L{4.6cm}@{}}
\toprule
\textbf{Признак} & \textbf{Свързване с libclang} & \textbf{Изпълнение на компилатора и разбор на JSON} \\
\midrule
Зависимост при построяване & заглавни файлове и библиотека на LLVM &
    няма, нужен е само компилатор за C++20 и CMake \\
Зависимост при изпълнение & същата библиотека, със съвпадаща версия &
    двоичният файл \code{clang++}, без изисквания към версията \\
Двоична съвместимост & обвързва проекта с ABI на конкретно издание на LLVM &
    няма ABI, границата е текст \\
Построяване на чужда машина & изисква пакетен управител или ръчна инсталация &
    работи веднага \\
Обхват без фронтенд & нищо не се построява &
    всичко освен извличането се построява и работи \\
Цена на прехода & извикване на функция &
    създаване на процес и разбор на текст, измерено в Раздел~\ref{sec:results} \\
Пълнота на данните & целият вътрешен модел на фронтенда &
    само това, което изводът съдържа, вж. по-долу \\
\bottomrule
\end{tabular}
\end{table}

Приета е четвъртата възможност. Решаващият довод е зависимостта при построяване. Проектът
трябва да се построява на чиста машина само с компилатор за C++20 и CMake, а свързването
с libclang би изисквало инсталация на LLVM със заглавните файлове за разработка още при
конфигурирането. При изпълнение на компилатора като външен процес фронтендът е вход на
инструмента, а не част от него: без него не работи само извличането, докато излъчването
на диаграми, пораждането на схема, пораждането на заготовки и проверката за разминаване
продължават да работят. Границата между двата процеса е текст, така че отпадат и
съвпадението на версии, и двоичната съвместимост.

Изборът има и цена, която е измерена, а не премълчана. Създаването на процес и разборът
на текста струват време, което при свързване с библиотека не се плаща. Освен това изводът
в JSON не съдържа всичко, което вътрешният модел на фронтенда съдържа: установено е при
проверка, че \term{аргументите на атрибутите} се губят, тоест \code{AnnotateAttr} се
извежда като възел с местоположение, но без низа, който е бил подаден на атрибута. Това
обстоятелство пряко определя как се записва съответствието клас към таблица и е разгледано
по-долу.

\subsection{Формат на извеждане на диаграмата}

Възможностите са пряко рисуване във векторен формат, извеждане в стандартизирания XMI
обмен за UML модели и извеждане в текстова нотация за диаграми. Прякото рисуване изисква
реализация на алгоритъм за подреждане на графа и измества фокуса на проекта от
моделирането към геометрията. XMI е точен и стандартен, но е нечетим за човек и е полезен
само ако насрещният инструмент вече го приема. Избрано е извеждане в две текстови
нотации, PlantUML~\cite{plantuml_guide} и Mermaid~\cite{mermaid_docs}, защото двете са
четими в необработен вид, версионират се като текст и се изобразяват от съществуващи
изобразители. Двете, а не една, защото различията между тях са проверка, че вътрешният
модел е независим от изходния формат. Тази проверка не е формална: тя откри реална
загуба, отчетена в Раздел~\ref{sec:results}.

\subsection{Запис на съответствието клас към таблица}

Разгледани са външен описателен файл, извеждане по съглашение за именуване и вграждане на
описанието в самия код. Външният файл се разминава с кода по същата причина, по която се
разминава и диаграмата. Съглашението за именуване е тихо и се чупи при първото изключение.
Избрано е описанието да живее в кода, което на нивото на модела се изобразява като UML
стереотип и етикетирани стойности. Така механизмът за разширяване на UML и механизмът за
анотиране на кода съвпадат по предназначение и преходът между тях е пряк.

Остава въпросът какъв да бъде носителят в кода. Естественият избор са атрибутите на
C++20~\cite{iso_cpp20}, но той е неприложим по причина, установена с проверка, а не
предположена: изводът на фронтенда в JSON записва атрибута като възел с диапазон в
изходния текст и \textbf{без неговите аргументи}, тоест низът в
\code{[[clang::annotate("table=books")]]} изобщо не достига до четеца. Приет е втори
носител, \code{static constexpr} член с инициализатор, чиято стойност се извежда изцяло,
включително низовият литерал. Изборът е стандартен C++20, не струва нищо по време на
изпълнение, чете се от човек, отварящ заглавния файл, и не зависи от разширение на
конкретен компилатор. Цената му е, че анотацията изглежда като член на класа; четецът
разпознава представката \code{bp\_} и я премества върху класификатора като стереотип или
етикетирана стойност, така че тя никога не се появява като атрибут в модела.

\subsection{Език, система за построяване и външни зависимости}

C++20 е приет за език на реализацията, защото инструментът работи върху C++ и
разработчикът му остава в същата среда, и защото по-ранна работа на същия автор върху
съпоставяне на класове и релационни схеми по време на компилация е решавала сродна задача
от другата ѝ страна~\cite{tsvetanov_bsc_orm}. За система за построяване е приет CMake.

Приетото ограничение за външните зависимости е по-строго от обичайното: \term{нула
задължителни външни библиотеки}. То не е вкус, а следствие от предназначението. Хранилището
е публично и непознат човек трябва да може да изпълни \code{cmake -S . -B build} и
\code{cmake --build build} и това да проработи от първия път. Построяване, което изисква
пакетен управител, е построяване, което никой не прави. Следствията са три и всяко е
съзнателно: за модулно тестване не е използвана външна рамка, а собствен изпълнител от
около сто реда, свързан с CTest чрез \code{add\_test}; за измерване на време не е
използвана външна библиотека, а \code{std::chrono::steady\_clock}; и за разбор на JSON е
написан собствен четец, чийто обхват е точно това, което изводът на фронтенда изисква.
Общо тези три решения добавят по-малко код, отколкото би добавила настройката на пакетен
управител, и премахват изцяло въпроса дали проектът ще се построи другаде.
